\chapter{Shortest Paths: Foundations}

In Chapter~6, we found minimum spanning trees that connect an entire network with
minimal total edge weight. But global economy does not guarantee personal convenience:
the path between two towns in an MST can be frustratingly circuitous. If your goal is
to travel from a specific starting point to every other vertex as quickly or cheaply as
possible, you do not want an MST; you want a shortest-path tree. This chapter lays the
foundations of shortest-path algorithms: the principle of edge relaxation, Dijkstra's
algorithm for nonnegative weights across various heap data structures, and the concept
of vertex potentials that bridges into advanced routing.

We work every algorithm by hand on a small graph. Tracing a five-vertex example until the
distance labels stop changing teaches more than any amount of staring at pseudocode.

\chapterline

\section{The Problem, Stated Carefully}
\label{ch07:problem}

Let $\Gph=(V,E)$ be a directed graph with $n=\abs{V}$ vertices and $m=\abs{E}$ edges. A
\emph{weight function} $w\colon E\to\R$ assigns a real number to each edge. (Undirected
graphs are the special case where every edge $\{u,v\}$ is replaced by the two directed
edges $(u,v)$ and $(v,u)$ of equal weight -- but beware, this trick is dangerous with
negative weights, as we will see.) The \emph{weight of a path}
$p = \langle v_0,v_1,\dots,v_k\rangle$ is the sum of its edge weights,
\[
   w(p) \;=\; \sum_{i=1}^{k} w(v_{i-1},v_i).
\]
The \emph{shortest-path distance} from $s$ to $t$ is
\[
   \dist(s,t) \;=\;
   \begin{cases}
     \min\set{ w(p) : p \text{ is a path from } s \text{ to } t}, & \text{if such a path exists},\\
     +\infty, & \text{if } t \text{ is unreachable from } s.
   \end{cases}
\]

\begin{definition}[Single-source shortest paths, SSSP]
\label{ch07:def-sssp}
Given $\Gph$, $w$, and a distinguished \emph{source} $s\in V$, compute $\dist(s,v)$ for
every $v\in V$, together with a \emph{shortest-path tree} that records one optimal route to
each reachable vertex.
\end{definition}

Two subtleties deserve attention before we touch an algorithm.

\begin{trapbox}
\textbf{Negative cycles make ``shortest'' meaningless.} Suppose a cycle $C$ reachable from
$s$ and able to reach $t$ has total weight $w(C)<0$. Then for any path from $s$ to $t$ we
can splice in another lap of $C$ to lower its weight, without bound. The infimum is
$-\infty$ and no shortest \emph{path} exists. Whenever negative weights are in play, ``what
is the shortest path?'' is only well posed once we have ruled out -- or detected -- a
negative cycle on the relevant part of the graph. Bellman--Ford
(Section~\ref{ch07:bellman-ford}) does exactly this detection.
\end{trapbox}

\begin{intuitionbox}
Even when distances are finite, why should a \emph{shortest path itself} be simple (no
repeated vertex)? Because if a path revisits a vertex, it contains a cycle. With no
negative cycle anywhere, every cycle has weight $\ge 0$, so deleting that cycle cannot
increase the weight -- the shortened walk is still a path and is no worse. Repeating the
argument, some shortest route uses at most $n-1$ edges. That bound -- at most $n-1$
edges -- is the key fact behind Bellman--Ford.
\end{intuitionbox}

\subsection{The optimal-substructure property}

Every shortest-path algorithm relies on the following fact.

\begin{lemma}[Subpaths of shortest paths are shortest]
\label{ch07:optsub}
Let $p=\langle v_0,\dots,v_k\rangle$ be a shortest path from $v_0$ to $v_k$, and let
$p_{ij}=\langle v_i,\dots,v_j\rangle$ be the contiguous subpath from $v_i$ to $v_j$, where
$0\le i\le j\le k$. Then $p_{ij}$ is a shortest path from $v_i$ to $v_j$.
\end{lemma}

\begin{proof}
Decompose $p$ into $v_0 \xrightarrow{p_{0i}} v_i \xrightarrow{p_{ij}} v_j
\xrightarrow{p_{jk}} v_k$, so that $w(p)=w(p_{0i})+w(p_{ij})+w(p_{jk})$. If some path
$p'_{ij}$ from $v_i$ to $v_j$ had $w(p'_{ij})<w(p_{ij})$, then replacing $p_{ij}$ by
$p'_{ij}$ yields a $v_0$-to-$v_k$ walk of weight strictly less than $w(p)$. With no
negative cycle this walk contains a path of no greater weight, contradicting the
optimality of $p$.
\end{proof}

Optimal substructure is what lets us speak of a shortest-path \emph{tree}: if the best
route to $v$ ends with the edge $(u,v)$, then the prefix ending at $u$ is itself a best
route to $u$. So we can store, for each vertex, only its predecessor on some optimal path,
and the union of these predecessor edges forms a tree rooted at $s$ that simultaneously
encodes a shortest path to every vertex.

\section{Relaxation: The One Move Every Algorithm Makes}
\label{ch07:relaxation}

Dijkstra, Bellman--Ford, and the DAG algorithm all repeat the same basic step,
\emph{relaxation}, and differ only in the order in which they apply it to edges.

We maintain for each vertex $v$ a tentative distance $d[v]$, an over-estimate of
$\dist(s,v)$ that we steadily push down, and a predecessor pointer $\pi[v]$. Initialization
sets $d[s]=0$, $d[v]=+\infty$ for $v\neq s$, and all $\pi[v]=\textsc{nil}$.

\begin{algorithm}[H]
\DontPrintSemicolon
\caption{\textsc{Relax}$(u,v,w)$ -- the atomic step}
\If{$d[u] + w(u,v) < d[v]$}{
  $d[v] \gets d[u] + w(u,v)$\;
  $\pi[v] \gets u$\;
}
\end{algorithm}

\begin{intuitionbox}
Relaxing the edge $(u,v)$ asks one question: ``Is going to $u$ and then taking this edge a
cheaper way to reach $v$ than anything I have found so far?'' If yes, we lower $d[v]$ and
remember that we arrived via $u$. The word ``relax'' is borrowed from physics: we let a
constraint that was held too tight (an over-estimate) loosen toward its true value.
\end{intuitionbox}

What makes relaxation safe is that no matter \emph{how} you order the steps, a label
never drops below the true shortest-path distance. Two invariants make this precise.

\begin{lemma}[Upper-bound invariant]
\label{ch07:ub-invariant}
At all times, $d[v]\ge \dist(s,v)$ for every $v$, and once $d[v]$ reaches $\dist(s,v)$ it
never changes again.
\end{lemma}

\begin{proof}
Initially $d[s]=0=\dist(s,s)$ (assuming no negative cycle through $s$) and $d[v]=\infty\ge
\dist(s,v)$ for the rest. Relaxation only ever lowers a label, and only to $d[u]+w(u,v)$. By
induction $d[u]\ge\dist(s,u)$, so $d[u]+w(u,v)\ge \dist(s,u)+w(u,v)\ge \dist(s,v)$, the last
step because $\dist(s,u)+w(u,v)$ is the weight of \emph{some} route to $v$ and the distance
is the minimum over all routes. Thus the new value still bounds $\dist(s,v)$ from above. The
value $\dist(s,v)$ is therefore a floor the label can never break through, so once reached
it is frozen.
\end{proof}

\begin{lemma}[Convergence / path-relaxation]
\label{ch07:path-relax}
Let $p=\langle s=v_0,v_1,\dots,v_k\rangle$ be a shortest path from $s$ to $v_k$. If at some
point the edges of $p$ are relaxed \emph{in order} $(v_0,v_1),(v_1,v_2),\dots,(v_{k-1},v_k)$
-- with arbitrary other relaxations interspersed -- then afterward $d[v_k]=\dist(s,v_k)$,
and it stays there.
\end{lemma}

\begin{proof}
By induction on $i$ we show $d[v_i]=\dist(s,v_i)$ holds right after $(v_{i-1},v_i)$ is
relaxed. Base: $d[v_0]=d[s]=0=\dist(s,s)$ from the start. Step: just before $(v_{i-1},v_i)$
is relaxed we already have $d[v_{i-1}]=\dist(s,v_{i-1})$, established earlier in the order
and frozen by Lemma~\ref{ch07:ub-invariant}. The relaxation tests $d[v_{i-1}]+w(v_{i-1},v_i)
= \dist(s,v_{i-1})+w(v_{i-1},v_i)=\dist(s,v_i)$, the last equality because $p$ is a shortest
path and its prefix to $v_i$ is shortest (Lemma~\ref{ch07:optsub}). So $d[v_i]$ is set to
$\dist(s,v_i)$ (or already equals it), then frozen.
\end{proof}

Together, the two lemmas say that every shortest-path algorithm is a \emph{schedule}
of relaxations chosen to guarantee that, for each vertex, the edges of \emph{some} shortest
path get relaxed in order. The algorithms differ only in how cleverly -- and how cheaply --
they arrange that.

\subsection{Optimality conditions}

When have we finished? The following ``no edge can be improved'' condition is the
certificate.

\begin{theorem}[Shortest-path optimality conditions]
\label{ch07:opt-conditions}
Suppose $\Gph$ contains no negative cycle reachable from $s$, and let $d[\cdot]$ be any
labeling produced by a sequence of relaxations from the standard start $d[s]=0$,
$d[v]=+\infty$ ($v\ne s$) --- or more generally any labeling satisfying $d[s]=0$ and the
upper-bound property $d[v]\ge\dist(s,v)$ for all $v\in V$. Then $d[\cdot]$ equals the true
distances $\dist(s,\cdot)$ \emph{if and only if} it is \emph{feasible}:
\[
   d[v] \;\le\; d[u] + w(u,v) \qquad \text{for every edge } (u,v)\in E . \tag{$\star$}
\]
\end{theorem}

\begin{proof}
($\Rightarrow$) If $d[v]=\dist(s,v)$ then $(\star)$ is just the triangle inequality for
distances: there is a route to $v$ via $u$ of weight $\dist(s,u)+w(u,v)$, and the distance
is no larger.

($\Leftarrow$) Suppose $(\star)$ holds. Take any path $p=\langle s,v_1,\dots,v_k=t\rangle$.
Summing $(\star)$ along its edges telescopes:
\[
   d[t]-d[s] = \sum_{i=1}^{k}\bigl(d[v_i]-d[v_{i-1}]\bigr)
            \le \sum_{i=1}^{k} w(v_{i-1},v_i) = w(p).
\]
Since $d[s]=0$ we get $d[t]\le w(p)$ for \emph{every} path, hence $d[t]\le\dist(s,t)$.
Combined with the upper bound $d[t]\ge\dist(s,t)$ (which holds for any labeling produced by
relaxation from a correct start, Lemma~\ref{ch07:ub-invariant}), equality follows.
\end{proof}

\begin{whybox}
Condition $(\star)$ is exactly ``no edge is still relaxable.'' So an algorithm is allowed
to stop the moment relaxation can make no further progress: feasibility \emph{is}
optimality. This is the same telescoping trick that powers linear-programming duality, and
when we meet \emph{potentials} (Section~\ref{ch07:potentials}) the quantity $d[v]$ will play
the role of a dual variable.
\end{whybox}

\section{A Running Example}
\label{ch07:running-graph}

Here is the graph we will trace by hand throughout the chapter. Source $s$, five vertices,
all weights nonnegative for now (we add a negative edge later).

\begin{center}
\begin{tikzpicture}[>=Stealth,line width=0.9pt,
   v/.style={circle,draw=nodeblue,fill=blue!6,minimum size=8mm,inner sep=0pt,font=\small},
   every edge/.style={draw=edgegray,->}]
  \node[v,fill=sourceblue!25,draw=sourceblue] (s) at (0,0)   {$s$};
  \node[v] (a) at (2.4,1.4) {$a$};
  \node[v] (b) at (2.4,-1.4){$b$};
  \node[v] (c) at (5.0,1.4) {$c$};
  \node[v] (t) at (5.0,-1.4){$t$};
  \path
    (s) edge node[above left,font=\scriptsize]{4} (a)
    (s) edge node[below left,font=\scriptsize]{2} (b)
    (a) edge node[above,font=\scriptsize]{3} (c)
    (b) edge node[left,font=\scriptsize]{1} (a)
    (b) edge node[below,font=\scriptsize]{5} (t)
    (c) edge node[right,font=\scriptsize]{2} (t)
    (a) edge[bend left=12] node[above right,font=\scriptsize]{6} (t);
\end{tikzpicture}
\end{center}

Edges: $s\!\to\!a\,(4)$, $s\!\to\!b\,(2)$, $b\!\to\!a\,(1)$, $a\!\to\!c\,(3)$,
$a\!\to\!t\,(6)$, $b\!\to\!t\,(5)$, $c\!\to\!t\,(2)$. The true distances, which we will
rediscover three different ways, are
\[
  \dist(s,s)=0,\quad \dist(s,b)=2,\quad \dist(s,a)=3,\quad \dist(s,c)=6,\quad \dist(s,t)=7.
\]
The optimal route to $a$ is $s\to b\to a$ (cost $2+1=3$, cheaper than the direct edge of
weight $4$). For $t$, compare the candidates $s\to b\to t = 2+5 = 7$ and
$s\to b\to a\to c\to t = 2+1+3+2 = 8$ and $s\to b\to a\to t = 2+1+6 = 9$; the winner is the
direct hop through $b$, so $\dist(s,t)=7$. Keep these numbers handy -- all three algorithms
must agree on them.

\section{Dijkstra's Algorithm}
\label{ch07:dijkstra}

\begin{historybox}
The foundations of shortest paths crystallized in the late 1950s. Edsger Dijkstra described his algorithm in 1959, having designed it in 1956 in about twenty minutes at a cafe table while shopping with his fianc{\'e}e, as a demonstration for the ARMAC computer. Dijkstra observed that when edge weights are nonnegative, vertices can be finalized in greedy order without ever retracting a distance label. Around the same time, Lester Ford Jr.\ (1956), Richard Bellman (1958), and Edward F.~Moore (1959) developed the iterative edge-relaxation algorithm for general weights, deriving it via Bellman's dynamic programming principle of optimality. Later refinements tackled the priority queue bottleneck: Fredman and Tarjan (1984) invented Fibonacci heaps to achieve $\Oh(m + n \log n)$, while integer priority queues by Thorup and others pushed bounded-weight shortest paths toward linear time.
\end{historybox}

\subsection{The greedy idea}

\begin{intuitionbox}
Imagine fire spreading from $s$ at unit speed along edges whose ``length'' is their weight.
The vertex that catches fire next is always the one with the smallest tentative distance
among those not yet burning -- because with nonnegative weights, no detour through a
\emph{farther} unburned vertex could ever reach it sooner. Dijkstra is exactly this:
repeatedly burn (finalize) the nearest unfinished vertex, then relax its outgoing edges to
update its neighbors.
\end{intuitionbox}

The diagram below shows Dijkstra's expanding wavefront at the moment $S=\{s,b,a\}$ has been settled and outgoing edges from $S$ relax the priority-queue frontier $Q=\{c,t\}$.

\begin{center}
\begin{tikzpicture}[
    >=Stealth,
    vnode/.style={circle, draw=sourceblue, fill=sourceblue!10, thick, minimum size=7.5mm, font=\small\bfseries},
    settlednode/.style={circle, draw=sourceblue!90!black, fill=sourceblue!25, very thick, minimum size=7.5mm, font=\small\bfseries},
    frontiernode/.style={circle, draw=quotientorange!90!black, fill=quotientorange!15, very thick, minimum size=7.5mm, font=\small\bfseries},
    unsettlednode/.style={circle, draw=gray!60, fill=gray!8, thick, minimum size=7.5mm, font=\small\bfseries},
    relaxedge/.style={draw=quotientorange!90!black, very thick, ->},
    treeedge/.style={draw=sourceblue!85!black, very thick, ->},
    unseenedge/.style={draw=gray!50, thick, ->, dashed},
    paneltitle/.style={font=\footnotesize\scshape, text=gray!85!black}
]
  % Wavefront background boundaries
  \draw[fill=sourceblue!8, draw=sourceblue!40, rounded corners=10pt, dashed] (-0.8, -1.7) rectangle (3.4, 2.4);
  \node[font=\footnotesize\bfseries, text=sourceblue!90!black] at (1.3, 2.1) {Settled Set $S$ ($\dist$ finalized)};

  \draw[fill=quotientorange!5, draw=quotientorange!40, rounded corners=10pt, dashed] (4.0, -1.7) rectangle (6.8, 2.4);
  \node[font=\footnotesize\bfseries, text=quotientorange!90!black] at (5.4, 2.1) {Frontier in $Q$};

  \draw[fill=gray!4, draw=gray!30, rounded corners=10pt, dashed] (7.4, -1.7) rectangle (9.8, 2.4);
  \node[font=\footnotesize\bfseries, text=gray!70!black] at (8.6, 2.1) {Unseen};

  % Settled vertices
  \node[settlednode] (s) at (0.2, 0.1) {$s$};
  \node[settlednode] (b) at (2.2, -0.6) {$b$};
  \node[settlednode] (a) at (2.2, 0.8) {$a$};

  \draw[treeedge] (s) -- node[midway, below=1pt, font=\scriptsize] {$2$} (b);
  \draw[treeedge] (b) -- node[midway, left=1pt, font=\scriptsize] {$1$} (a);

  % Expanding Wavefront labels
  \node[font=\scriptsize\bfseries, text=sourceblue!90!black, above=2pt of a] {$d[a]=3$};
  \node[font=\scriptsize\bfseries, text=sourceblue!90!black, below=2pt of b] {$d[b]=2$};
  \node[font=\scriptsize\bfseries, text=sourceblue!90!black, left=2pt of s] {$d[s]=0$};

  % Frontier vertices
  \node[frontiernode] (c) at (5.4, 0.8) {$c$};
  \node[frontiernode] (t) at (5.4, -0.6) {$t$};

  \node[font=\scriptsize\bfseries, text=quotientorange!95!black, above=2pt of c] {$d[c]=6$};
  \node[font=\scriptsize\bfseries, text=quotientorange!95!black, below=2pt of t] {$d[t]=7$};

  % Relaxing edges across wavefront
  \draw[relaxedge] (a) -- node[midway, above=1pt, font=\scriptsize\bfseries, text=quotientorange!95!black] {relax $+3$} (c);
  \draw[relaxedge] (b) -- node[midway, below=1pt, font=\scriptsize\bfseries, text=quotientorange!95!black] {relax $+5$} (t);
  \draw[unseenedge] (a) -- node[pos=0.35, below left=1pt, font=\scriptsize] {$+6$} (t);

  % Unseen vertex
  \node[unsettlednode] (z) at (8.6, 0.1) {$z$};
  \node[font=\scriptsize, text=gray!70!black, above=2pt of z] {$d[z]=\infty$};
  \draw[unseenedge] (c) -- node[midway, above=1pt, font=\scriptsize] {$4$} (z);

  % Wavefront legend
  \node[anchor=center, font=\scriptsize] at (4.5, -2.2) {
    \textbf{Expanding Wavefront:}\quad
    \tikz[baseline=-0.5ex]{\node[circle,draw=sourceblue,fill=sourceblue!25,inner sep=2pt,minimum size=3.5mm]{};} Settled \qquad
    \tikz[baseline=-0.5ex]{\node[circle,draw=quotientorange,fill=quotientorange!15,inner sep=2pt,minimum size=3.5mm]{};} Priority Queue Frontier \qquad
    \tikz[baseline=-0.5ex]{\draw[relaxedge] (0,0) -- (0.5,0);} Active Relaxation
  };
\end{tikzpicture}
\end{center}

Maintain a set $S$ of \emph{settled} vertices whose distance is final, and a priority queue
$Q$ of the rest keyed by $d[\cdot]$. Each round, extract the minimum-key vertex $u$ from
$Q$, move it to $S$, and relax every edge leaving $u$.

\begin{algorithm}[H]
\DontPrintSemicolon
\caption{\textsc{Dijkstra}$(\Gph,w,s)$ -- nonnegative weights}
\KwIn{Directed graph $\Gph=(V,E)$, weights $w\ge 0$, source $s$.}
\KwOut{$d[v]=\dist(s,v)$ and predecessors $\pi[\cdot]$.}
\ForEach{$v\in V$}{ $d[v]\gets\infty$;\quad $\pi[v]\gets\textsc{nil}$ }
$d[s]\gets 0$\;
$S\gets\emptyset$;\quad build min-priority-queue $Q$ over $V$ keyed by $d$\;
\While{$Q\neq\emptyset$}{
  $u\gets\textsc{Extract-Min}(Q)$\;
  $S\gets S\cup\{u\}$\;
  \ForEach{edge $(u,v)\in E$}{
     \If{$d[u]+w(u,v) < d[v]$}{
        $d[v]\gets d[u]+w(u,v)$;\quad $\pi[v]\gets u$\;
        \textsc{Decrease-Key}$(Q,v,d[v])$\;
     }
  }
}
\end{algorithm}

\subsection{Correctness}

\begin{whybox}
The whole proof rests on a single claim: \emph{when a vertex $u$ is extracted, $d[u]$
already equals $\dist(s,u)$.} If that holds, then relaxing $u$'s edges relaxes, in order,
the last edge of every shortest path that ends by passing through $u$, and
Lemma~\ref{ch07:path-relax} finishes the argument. Proving the claim is where nonnegativity is
needed.
\end{whybox}

\begin{theorem}[Dijkstra is correct]
\label{ch07:dijkstra-correct}
If $w(e)\ge 0$ for all $e$, then on termination $d[v]=\dist(s,v)$ for every $v\in V$.
\end{theorem}

\begin{proof}
We prove the invariant: \emph{each vertex $u$ has $d[u]=\dist(s,u)$ at the instant it is
extracted from $Q$.} By Lemma~\ref{ch07:ub-invariant}, $d[u]\ge\dist(s,u)$ always, so we
need only rule out $d[u]>\dist(s,u)$ at extraction time.

Suppose, for contradiction, that $u$ is the \emph{first} vertex extracted with
$d[u]>\dist(s,u)$. Then $u\neq s$ and $u$ is reachable, so a shortest path
$p\colon s\rightsquigarrow u$ exists with $w(p)=\dist(s,u)<\infty$. Walk along $p$ from $s$;
since $s\in S$ already (it was extracted first, correctly) and $u\notin S$ at this moment,
$p$ has a first edge $(x,y)$ that leaves $S$, i.e. $x\in S$ and $y\notin S$ (possibly
$y=u$).

Because $x$ was extracted before $u$ and was \emph{not} a counterexample (we took $u$ to be
the first), we had $d[x]=\dist(s,x)$ when $x$ was settled, and the edge $(x,y)$ was relaxed
at that time. Hence
\[
   d[y] \;\le\; d[x]+w(x,y) \;=\; \dist(s,x)+w(x,y) \;=\; \dist(s,y),
\]
the final equality because $(x,y)$ lies on a shortest path to $u$ and therefore its prefix
to $y$ is shortest (Lemma~\ref{ch07:optsub}). Combined with the upper bound, $d[y]=
\dist(s,y)$.

Now use nonnegativity: $y$ lies on $p$ at or before $u$, and the segment of $p$ from $y$ to
$u$ has weight $\ge 0$, so $\dist(s,y)\le\dist(s,u)$. Putting the pieces together,
\[
   d[y]=\dist(s,y)\le\dist(s,u)<d[u].
\]
But $u$ was chosen as the minimum-key vertex over $Q$, and $y\in Q$ at extraction time, so
$d[u]\le d[y]$ -- a contradiction. Hence no such first counterexample exists; every vertex
is extracted with its correct distance. After extraction its edges are relaxed, so every
shortest path's edges are eventually relaxed in order, and
Lemma~\ref{ch07:path-relax} gives $d[v]=\dist(s,v)$ for all $v$.
\end{proof}

\begin{trapbox}
\textbf{Dijkstra and negative edges do not mix.} The line ``the segment from $y$ to $u$ has
weight $\ge 0$, so $\dist(s,y)\le\dist(s,u)$'' is the \emph{only} place nonnegativity is
used -- and it is fatal without it. Concretely, take $s\to a$ of weight $1$, $s\to b$ of
weight $2$, and $b\to a$ of weight $-2$. Dijkstra extracts $a$ first with $d[a]=1$, settles
it, and never revisits; but $\dist(s,a)=2-2=0$. The greedy commitment is wrong. Do not
``fix'' this by re-inserting settled vertices and hoping -- that degenerates to exponential
time on adversarial graphs. Use Bellman--Ford, or reweight with potentials.
\end{trapbox}

\subsection{Tracing Dijkstra by hand}

\begin{labbox}
Run Dijkstra on the graph of Section~\ref{ch07:running-graph}. Columns are the labels
$d[\cdot]$; ``--'' means $\infty$. The settled vertex of each round is starred.

\medskip
\begin{center}
\begin{tabular}{llccccc}
\toprule
Round & Extract & $d[s]$ & $d[a]$ & $d[b]$ & $d[c]$ & $d[t]$ \\
\midrule
init  & --        & $0$ & -- & -- & -- & -- \\
1 & $s^\star$ & $0$ & $4$ & $2$ & -- & -- \\
2 & $b^\star$ & $0$ & $3$ & $2$ & -- & $7$ \\
3 & $a^\star$ & $0$ & $3$ & $2$ & $6$ & $7$ \\
4 & $c^\star$ & $0$ & $3$ & $2$ & $6$ & $7$ \\
5 & $t^\star$ & $0$ & $3$ & $2$ & $6$ & $7$ \\
\bottomrule
\end{tabular}
\end{center}

\medskip
Walk it: Round 1 settles $s$, relaxes $s\to a$ (to $4$) and $s\to b$ (to $2$). Round 2 the
queue minimum is $b\,(2)$, not $a\,(4)$ -- settle $b$; relax $b\to a$ giving $2+1=3<4$ (down
to $3$, predecessor now $b$) and $b\to t$ giving $2+5=7$. Round 3 settle $a\,(3)$; relax
$a\to c$ to $3+3=6$ and $a\to t$ to $3+6=9$, which is \emph{not} less than the current $7$,
so $d[t]$ stays $7$. Round 4 settle $c\,(6)$; relax $c\to t$ to $6+2=8\not<7$, no change.
Round 5 settle $t\,(7)$. Final tree edges: $\pi[a]=b,\ \pi[b]=s,\ \pi[c]=a,\ \pi[t]=b$.
Notice $a$'s label dropped from $4$ to $3$ \emph{before} $a$ was settled -- that is the
decrease-key in action.
\end{labbox}

The predecessor pointers collect into the shortest-path tree below: each vertex is labeled
with its final distance, and the bold arrows are the $\pi$ edges $s\to b\to a$, $a\to c$,
and $b\to t$. The original edges $s\to a$, $a\to t$, and $c\to t$ are dashed -- they were relaxed but
lost, never becoming part of any optimal route.

\begin{figure}[H]
\centering
\begin{tikzpicture}[>=Stealth,line width=0.9pt,
   v/.style={circle,draw=nodeblue,fill=blue!6,minimum size=9mm,inner sep=0pt,font=\small},
   tree/.style={draw=sourceblue,->,line width=1.4pt},
   lost/.style={draw=edgegray,->,dashed}]
  \node[v,fill=sourceblue!25,draw=sourceblue] (s) at (0,0)   {$s$};
  \node[v] (a) at (2.4,1.4) {$a$};
  \node[v] (b) at (2.4,-1.4){$b$};
  \node[v] (c) at (5.0,1.4) {$c$};
  \node[v] (t) at (5.0,-1.4){$t$};
  \node[font=\scriptsize,above=0.5pt of s] {$0$};
  \node[font=\scriptsize,above=0.5pt of a] {$3$};
  \node[font=\scriptsize,below=0.5pt of b] {$2$};
  \node[font=\scriptsize,above=0.5pt of c] {$6$};
  \node[font=\scriptsize,below=0.5pt of t] {$7$};
  \path
    (s) edge[lost] node[above left,font=\scriptsize]{4} (a)
    (s) edge[tree] node[below left,font=\scriptsize]{2} (b)
    (a) edge[tree] node[above,font=\scriptsize]{3} (c)
    (b) edge[tree] node[left,font=\scriptsize]{1} (a)
    (b) edge[tree] node[below,font=\scriptsize]{5} (t)
    (c) edge[lost] node[right,font=\scriptsize]{2} (t)
    (a) edge[lost,bend left=12] node[above right,font=\scriptsize]{6} (t);
\end{tikzpicture}
\caption{Shortest-path tree from Dijkstra on the running graph. Bold blue arrows are the
final $\pi$ edges; dashed gray arrows were relaxed but never won. Vertex labels are final
distances.}
\label{ch07:spt-figure}
\end{figure}

\subsection{Heap variants and complexity}
\label{ch07:dijkstra-heaps}

Dijkstra performs $n$ \textsc{Extract-Min} operations, at most $m$ \textsc{Decrease-Key}
operations, and $n$ inserts. Its running time is therefore
\[
   T = n\cdot T_{\text{insert}} + n\cdot T_{\text{ext}} + m\cdot T_{\text{dec}},
\]
and the priority queue we plug in decides everything. A standard binary min-heap stores
elements in a complete binary tree of height $\lfloor \log_2 n \rfloor$ satisfying the heap property
(each parent has a smaller key than its children). Inserting a new vertex or decreasing a key sifts
up along a single branch in $\Oh(\log n)$ steps, while extracting the minimum replaces the root with
the last leaf and sifts down in $\Oh(\log n)$ steps.

\begin{center}
\begin{tabular}{lcccl}
\toprule
Priority queue & Insert & Extract-Min & Decrease-Key & Total \\
\midrule
Array / unsorted    & $\Oh(1)$ & $\Oh(n)$ & $\Oh(1)$ & $\Oh(n^2 + m)$ \\
Binary heap         & $\Oh(\log n)$ & $\Oh(\log n)$ & $\Oh(\log n)$ & $\Oh((n+m)\log n)$ \\
$d$-ary heap, $d{=}m/n$ & $\Oh(\log_d n)$ & $\Oh(d\log_d n)$ & $\Oh(\log_d n)$ & $\Oh(m\log_{m/n} n)$ \\
Fibonacci heap      & $\Oh(1)$ & $\Oh(\log n)^\dagger$ & $\Oh(1)^\dagger$ & $\Oh(m + n\log n)$ \\
\bottomrule
\end{tabular}
\end{center}

\noindent{\footnotesize $\dagger$ amortized.}

\begin{intuitionbox}
The trade is between extraction and decrease-key. A binary heap makes both
$\Oh(\log n)$, fine when $m=\Th(n)$. But dense graphs have $m=\Th(n^2)$ decrease-keys,
and paying $\log n$ for each hurts. Two escapes. A \emph{$d$-ary heap} (each node has $d$
children) makes the tree shallower: height $\log_d n$ speeds decrease-key, at the cost of
scanning $d$ children per extract. Balancing $d\approx m/n$ gives $\Oh(m\log_{m/n}n)$,
which is $\Oh(m)$ for sufficiently dense graphs. A \emph{Fibonacci heap} makes
decrease-key $\Oh(1)$ amortized outright, yielding the classic $\Oh(m+n\log n)$ bound that
is optimal in the comparison model for this operation mix.
\end{intuitionbox}

Although Fibonacci heaps achieve the $\Oh(m+n\log n)$ bound by reducing decrease-key to
$\Oh(1)$ amortized time, in practice they carry significant pointer overhead and poor cache
locality, so engineers typically prefer binary or $4$-ary heaps, or a lazy-deletion binary
heap that re-inserts updated keys and discards stale entries.

\begin{connectionbox}
The family of heap structures -- binary, $d$-ary, Fibonacci -- is developed for its own sake in the
companion \emph{Algorithms From the Ground Up} (priority-queue chapter) and in the
data-structures appendix of this book. Dijkstra's algorithm is the primary workload motivating all three designs. The exchange
argument that proves Dijkstra correct is the same greedy-with-a-cut argument that proves
Prim's MST algorithm correct (the minimum-spanning-tree chapter); indeed Prim's algorithm
\emph{is} Dijkstra with ``$w(u,v)$'' in place of ``$d[u]+w(u,v)$'' as the key.
\end{connectionbox}

\begin{complexitybox}
Dijkstra, SSSP, nonnegative weights:
\begin{itemize}
  \item Binary heap: $\Oh((n+m)\log n)$ time, $\Oh(n)$ extra space. The default.
  \item Fibonacci heap: $\Oh(m+n\log n)$ time -- best comparison-based bound.
  \item Array: $\Oh(n^2)$ time -- best when $m=\Th(n^2)$ (dense) and simplest to code.
\end{itemize}
\end{complexitybox}

\subsection{The lazy-deletion implementation}

In practice many languages lack a built-in decrease-key. The standard workaround is to push
$(d[v],v)$ afresh every time $d[v]$ drops and to discard stale pairs on pop.

\begin{algorithm}[H]
\DontPrintSemicolon
\caption{\textsc{Dijkstra-Lazy}$(\Gph,w,s)$ -- binary heap, no decrease-key}
\ForEach{$v\in V$}{ $d[v]\gets\infty$ } \quad $d[s]\gets 0$\;
$Q\gets$ empty min-heap; push $(0,s)$\;
\While{$Q\neq\emptyset$}{
  $(\delta,u)\gets\textsc{Pop}(Q)$\;
  \lIf{$\delta>d[u]$}{\textbf{continue} \tcp*[f]{stale entry, skip}}
  \ForEach{$(u,v)\in E$}{
     \If{$d[u]+w(u,v)<d[v]$}{
        $d[v]\gets d[u]+w(u,v)$;\ $\pi[v]\gets u$;\ push $(d[v],v)$\;
     }
  }
}
\end{algorithm}

The heap may hold up to $m$ entries, so each operation is $\Oh(\log m)=\Oh(\log n)$ (since
$m\le n^2$), and the total is $\Oh(m\log n)$. The stale-skip line is what keeps it correct:
a vertex may be popped several times, but only its first (smallest) pop does real work.

\section{Bellman--Ford and Negative Cycles}
\label{ch07:bellman-ford}

Dijkstra's speed comes from settling each vertex once -- and that is exactly what breaks
under negative edges. Bellman--Ford pays for generality by giving up the greedy order: it
simply relaxes \emph{every} edge, over and over, $n-1$ times.

\begin{intuitionbox}
Recall (Section~\ref{ch07:problem}) that with no negative cycle, some shortest path uses at
most $n-1$ edges. Lemma~\ref{ch07:path-relax} says that if a shortest path's edges get
relaxed \emph{in order}, the endpoint converges. We do not know the order -- so we relax all
$m$ edges, then all $m$ again, $n-1$ rounds total. After round $i$, every shortest path of
$\le i$ edges has converged: its edges appear in order across the first $i$ sweeps,
because each sweep touches every edge. So after $n-1$ rounds, all shortest paths have
converged.
\end{intuitionbox}

\begin{algorithm}[H]
\DontPrintSemicolon
\caption{\textsc{Bellman-Ford}$(\Gph,w,s)$ -- arbitrary weights}
\ForEach{$v\in V$}{ $d[v]\gets\infty$;\ $\pi[v]\gets\textsc{nil}$ } \quad $d[s]\gets 0$\;
\For{$i\gets 1$ \KwTo $n-1$}{
   \ForEach{edge $(u,v)\in E$}{
      \If{$d[u]+w(u,v) < d[v]$}{ $d[v]\gets d[u]+w(u,v)$;\ $\pi[v]\gets u$ }
   }
}
\tcp{Negative-cycle check: one more sweep must change nothing.}
\ForEach{edge $(u,v)\in E$}{
   \lIf{$d[u]+w(u,v) < d[v]$}{\Return \textsc{NegativeCycleDetected}}
}
\Return $(d,\pi)$\;
\end{algorithm}

\begin{theorem}[Bellman--Ford correctness]
\label{ch07:bf-correct}
If $\Gph$ has no negative cycle reachable from $s$, then after the $n-1$ relaxation rounds,
$d[v]=\dist(s,v)$ for every $v$, and the final check reports no cycle. If a reachable
negative cycle \emph{does} exist, the final check returns \textsc{NegativeCycleDetected}.
\end{theorem}

\begin{proof}
\emph{No negative cycle.} Fix a vertex $v$ and a shortest path
$p=\langle s=v_0,\dots,v_k=v\rangle$ to it with $k\le n-1$ edges
(such a simple shortest path exists, Section~\ref{ch07:problem}). We claim by induction that
after round $i$, $d[v_i]=\dist(s,v_i)$. Each round relaxes \emph{all} edges, in particular
$(v_{i-1},v_i)$. After round $1$ the edge $(v_0,v_1)$ has been relaxed with
$d[v_0]=0$ already correct, fixing $d[v_1]$. Assuming $d[v_{i-1}]$ correct after round
$i-1$, round $i$ relaxes $(v_{i-1},v_i)$ and fixes $d[v_i]$ by
Lemma~\ref{ch07:path-relax}. As $k\le n-1$, after $n-1$ rounds $d[v_k]=\dist(s,v)$. Then
$(\star)$ holds for every edge, so the final sweep changes nothing.

\emph{Negative cycle exists.} Let $C=\langle v_0,\dots,v_\ell=v_0\rangle$ be a reachable
cycle with $\sum_{i} w(v_{i-1},v_i)<0$. Suppose the final check found \emph{no} improvable
edge; then $d[v_i]\le d[v_{i-1}]+w(v_{i-1},v_i)$ for all edges of $C$ and all $d[v_i]$ are
finite (reachability). Summing around the cycle,
\[
   \sum_{i=1}^{\ell} d[v_i] \;\le\; \sum_{i=1}^{\ell} d[v_{i-1}] + \sum_{i=1}^{\ell}
   w(v_{i-1},v_i).
\]
The two distance sums are equal (the cycle returns to its start), so they cancel, leaving
$0\le w(C)<0$, a contradiction. Hence some edge of $C$ is still improvable and the check
fires.
\end{proof}

\begin{whybox}
The negative-cycle test is the telescoping argument of
Theorem~\ref{ch07:opt-conditions} run \emph{around a loop}: feasibility $(\star)$ summed
over a cycle yields $0\le w(C)$, so feasibility is achievable exactly when every reachable
cycle is nonnegative. A surviving improvable edge after $n-1$ rounds is thus a
\emph{certificate} of a negative cycle. To extract the cycle itself, follow $\pi$ pointers
back $n$ steps from the improvable edge's head; you are guaranteed to land inside the
cycle, and tracing $\pi$ until a vertex repeats prints it.
\end{whybox}

\subsection{Tracing Bellman--Ford with a negative edge}

\begin{labbox}
Modify the running graph: change $b\to a$ from $+1$ to $-1$. The graph stays acyclic, so no
negative cycle appears. Now $\dist(s,a)=2+(-1)=1$, $\dist(s,c)=1+3=4$, and
$\dist(s,t)=\min(2+5,\,1+6,\,4+2)=6$. Process edges in the fixed order
$sa,sb,ba,ac,at,bt,ct$.

\medskip
\begin{center}
\begin{tabular}{lccccc}
\toprule
After round & $d[s]$ & $d[a]$ & $d[b]$ & $d[c]$ & $d[t]$ \\
\midrule
init & $0$ & -- & -- & -- & -- \\
1 & $0$ & $1$ & $2$ & $4$ & $6$ \\
2 & $0$ & $1$ & $2$ & $4$ & $6$ \\
\bottomrule
\end{tabular}
\end{center}

\medskip
Round 1, edge by edge: $sa$ sets $d[a]=4$; $sb$ sets $d[b]=2$; $ba$ sets
$d[a]=2-1=1$; $ac$ sets $d[c]=1+3=4$; $at$ sets $d[t]=1+6=7$; $bt$ tries $2+5=7$
(no change); $ct$ sets $d[t]=4+2=6$. Already correct after one round, because this edge
order happens to follow the shortest paths. Round 2 changes nothing -- and a third sweep
(the check) also changes nothing, certifying no negative cycle. The lucky single-round
convergence is order-dependent; reverse the edge list and you would need more rounds.
\end{labbox}

\begin{trapbox}
\textbf{An undirected negative edge immediately creates a negative two-cycle.} If you model an
undirected edge $\{u,v\}$ of weight $-3$ as the pair $(u,v),(v,u)$ both of weight $-3$,
you have created a $2$-cycle of weight $-6$. Bellman--Ford will (correctly) report a
negative cycle. Undirected shortest paths with negative edges is genuinely harder -- it
reduces to minimum-weight perfect matching, not to this chapter's machinery.
\end{trapbox}

\subsection{Complexity and the early-exit refinement}

Each of the $n-1$ rounds touches all $m$ edges, so Bellman--Ford runs in $\Oh(nm)$ time and
$\Oh(n)$ space. A useful refinement: if an entire round relaxes \emph{no} edge, the labels
are already feasible and we can stop early. The worst case stays $\Oh(nm)$, but on typical
inputs it finishes in a few rounds.

\begin{complexitybox}
Bellman--Ford, SSSP, arbitrary weights: $\Oh(nm)$ time, $\Oh(n)$ space, detects negative
cycles. The \emph{SPFA} heuristic (queue only the vertices whose label changed) has the same
$\Oh(nm)$ worst case but is often much faster; do not rely on it against adversarial graphs.
\end{complexitybox}

Viewed as a dynamic program, Bellman--Ford computes shortest-path distances stratified by hop count. By defining
$\dist^{(i)}(s,v)$ as the shortest distance using at most $i$ edges, the DP recurrence
\[
  \dist^{(i)}(s,v)=\min\Bigl(\dist^{(i-1)}(s,v),\;
     \min_{(u,v)\in E}\bigl[\dist^{(i-1)}(s,u)+w(u,v)\bigr]\Bigr)
\]
relaxes all edges in rounds, forming the direct ancestor of the Floyd--Warshall equations.

\section{Shortest Paths in a DAG}
\label{ch07:dag}

When the graph is a directed acyclic graph, negative weights cause no trouble (no cycles, so
certainly no negative ones), and we can do better than Bellman--Ford and even avoid the
priority queue of Dijkstra. The key is a \emph{topological order}: a linear arrangement of
the vertices such that every edge points forward.

\begin{intuitionbox}
In topological order, by the time we process a vertex $u$ every edge \emph{into} $u$ has
already been relaxed -- because those edges come from earlier vertices. So $d[u]$ is final
the instant we reach $u$, with no need to find the global minimum. One left-to-right sweep,
relaxing each vertex's outgoing edges as we pass it, settles the whole graph. It is
Bellman--Ford's $n-1$ rounds collapsed into a single perfectly-ordered pass.
\end{intuitionbox}

Placing the vertices of the running graph along a line in topological order $s,b,a,c,t$ makes every edge point strictly from left to right, as shown below.

\begin{center}
\begin{tikzpicture}[>=Stealth,line width=0.9pt,
   v/.style={circle,draw=nodeblue,fill=blue!6,minimum size=8mm,inner sep=0pt,font=\small},
   every edge/.style={draw=edgegray,->}]
  \node[v,fill=sourceblue!25,draw=sourceblue] (s) at (0,0) {$s$};
  \node[v] (b) at (1.8,0) {$b$};
  \node[v] (a) at (3.6,0) {$a$};
  \node[v] (c) at (5.4,0) {$c$};
  \node[v] (t) at (7.2,0) {$t$};
  \path
    (s) edge[bend left=30] node[above,font=\scriptsize]{4} (a)
    (s) edge node[above,font=\scriptsize]{2} (b)
    (b) edge node[above,font=\scriptsize]{1} (a)
    (a) edge node[above,font=\scriptsize]{3} (c)
    (c) edge node[above,font=\scriptsize]{2} (t)
    (b) edge[bend right=34] node[below,font=\scriptsize]{5} (t)
    (a) edge[bend left=30] node[above,font=\scriptsize]{6} (t);
\end{tikzpicture}
\end{center}

\begin{algorithm}[H]
\DontPrintSemicolon
\caption{\textsc{DAG-Shortest-Paths}$(\Gph,w,s)$}
$L\gets$ \textsc{TopologicalSort}$(\Gph)$ \tcp*{e.g.\ by DFS finish times}
\ForEach{$v\in V$}{ $d[v]\gets\infty$;\ $\pi[v]\gets\textsc{nil}$ } \quad $d[s]\gets 0$\;
\ForEach{$u\in L$ \textbf{\emph{in topological order}}}{
   \ForEach{$(u,v)\in E$}{
      \If{$d[u]+w(u,v) < d[v]$}{ $d[v]\gets d[u]+w(u,v)$;\ $\pi[v]\gets u$ }
   }
}
\Return $(d,\pi)$\;
\end{algorithm}

\begin{theorem}[DAG correctness]
\label{ch07:dag-correct}
On a DAG, \textsc{DAG-Shortest-Paths} computes $d[v]=\dist(s,v)$ for all $v$, even with
negative edge weights, in $\Oh(n+m)$ time.
\end{theorem}

\begin{proof}
Process vertices in topological order $u_1,u_2,\dots,u_n$. We show by induction on the
position $j$ that when the sweep finishes processing $u_j$, $d[u_j]=\dist(s,u_j)$. Any
shortest path to $u_j$ arrives via some last edge $(u_i,u_j)$ with $i<j$ (topological order
forces predecessors earlier), or $u_j=s$. If $u_j=s$ then $d[s]=0$. Otherwise, by induction
$d[u_i]=\dist(s,u_i)$ was settled when $u_i$ was processed, at which moment the edge
$(u_i,u_j)$ was relaxed, giving $d[u_j]\le \dist(s,u_i)+w(u_i,u_j)=\dist(s,u_j)$. With the
upper bound this is equality. Topological sort is $\Oh(n+m)$ and the sweep relaxes each edge
once, so the total is $\Oh(n+m)$.
\end{proof}

\begin{labbox}
The running graph is acyclic; a topological order is $s,b,a,c,t$ (every edge points right in
that list: $sa,sb,ba,ac,at,bt,ct$ all go forward). Sweep: process $s$ (relax $sa\!\to\!4$,
$sb\!\to\!2$); process $b$ (relax $ba\!\to\!3$, $bt\!\to\!7$); process $a$
(relax $ac\!\to\!6$, $at\!\to\!9\!\not<\!7$); process $c$
(relax $ct\!\to\!8\!\not<\!7$); process $t$ (no out-edges). Final
$d=(0,3,2,6,7)$ for $(s,a,b,c,t)$ -- identical to Dijkstra, obtained in one pass with no
heap. Replace any weight by a negative number and the same single pass still works.
\end{labbox}

\begin{applicationbox}
\textbf{Critical-path / PERT scheduling.} Reverse the sign and take \emph{longest} paths in
a DAG of tasks (each edge a precedence constraint carrying a duration); the longest
source-to-sink path is the project's critical path, the minimum possible completion time.
Longest path in a general graph is NP-hard, but in a DAG it is this same linear sweep with
$\max$ in place of $\min$. This is why project-management tools and build systems
(\texttt{make}, Bazel) lean on topological order.
\end{applicationbox}

\section{Potentials and Reduced Costs}
\label{ch07:potentials}

We close with an idea that turns negative-weight graphs into nonnegative ones \emph{without
changing which paths are shortest}, so that we can run fast Dijkstra where slow Bellman--Ford
seemed mandatory. Johnson's all-pairs algorithm and the successive-shortest-path
min-cost-flow methods both rely on it.

\begin{definition}[Potential and reduced cost]
\label{ch07:def-potential}
A \emph{potential} is any function $h\colon V\to\R$. It induces \emph{reduced edge weights}
\[
   w_h(u,v) \;=\; w(u,v) + h(u) - h(v).
\]
\end{definition}

\begin{intuitionbox}
Think of $h(v)$ as an altitude assigned to each vertex. Reducing the weight of $(u,v)$ by
$h(v)-h(u)$ charges you for climbing and credits you for descending. Along any path the
altitude terms telescope, so the \emph{relative} ranking of paths between two
fixed endpoints is untouched. We are free to choose altitudes that make every edge's reduced
weight nonnegative.
\end{intuitionbox}

\begin{lemma}[Reduced costs preserve shortest paths]
\label{ch07:reduced-preserves}
For any path $p$ from $u$ to $v$, $\;w_h(p) = w(p) + h(u) - h(v)$. Consequently $p$ is a
shortest $u$-$v$ path under $w$ if and only if it is shortest under $w_h$, and any cycle has
the same weight under $w$ and $w_h$.
\end{lemma}

\begin{proof}
Sum the definition along $p=\langle u=x_0,\dots,x_k=v\rangle$:
\[
   w_h(p)=\sum_{i=1}^{k}\bigl(w(x_{i-1},x_i)+h(x_{i-1})-h(x_i)\bigr)
         = w(p) + h(x_0)-h(x_k) = w(p)+h(u)-h(v),
\]
the middle terms telescoping. Since $h(u)-h(v)$ is a constant shared by all $u$-$v$ paths,
it shifts every path's weight identically and cannot change which is least. For a cycle,
$u=v$, so $h(u)-h(v)=0$ and the weight is unchanged.
\end{proof}

\begin{theorem}[Feasible potentials make weights nonnegative]
\label{ch07:feasible-potential}
The reduced weights $w_h$ are all nonnegative if and only if $h$ is a \emph{feasible
potential}: $h(v)\le h(u)+w(u,v)$ for every edge, i.e. $h$ satisfies the optimality
condition $(\star)$. Such an $h$ exists if and only if $\Gph$ has no negative cycle.
\end{theorem}

\begin{proof}
$w_h(u,v)=w(u,v)+h(u)-h(v)\ge 0 \iff h(v)\le h(u)+w(u,v)$, which is exactly $(\star)$.
By Theorem~\ref{ch07:opt-conditions}, choosing $h(v)=\dist(s,v)$ from any source $s$ that
reaches all of $V$ gives a feasible potential. To reach every vertex, add an artificial
source $q$ with a zero-weight edge $q\to v$ to all $v$; then $h(v)=\dist(q,v)$ is finite and
feasible exactly when no negative cycle exists (otherwise some $\dist(q,v)=-\infty$, and
Lemma~\ref{ch07:reduced-preserves} shows a negative cycle keeps its negative weight under
\emph{any} $h$, so no feasible $h$ could exist).
\end{proof}

\begin{whybox}
This is the bridge to linear-programming duality. The shortest-path problem is a linear
program; its dual variables are precisely the potentials $h(v)$, and ``feasible potential''
is dual feasibility. The condition $(\star)$ -- which we first met as ``no edge is
relaxable'' -- is dual feasibility, the optimality $d[v]=\dist(s,v)$ is primal--dual
complementary slackness, and the negative-cycle obstruction is dual infeasibility. The same
three words (feasible, reduced, slack) reappear verbatim in the min-cost-flow chapter.
\end{whybox}

\begin{labbox}
Take the running graph with $b\to a$ set to $-1$ (the negative variant). Use the
potential $h(v)=\dist(s,v)=(0,1,2,4,6)$ for $(s,a,b,c,t)$. Reduced weights:
\[
\begin{aligned}
 w_h(s,a)&=4+0-1=3, & w_h(s,b)&=2+0-2=0, & w_h(b,a)&=-1+2-1=0,\\
 w_h(a,c)&=3+1-4=0, & w_h(a,t)&=6+1-6=1, & w_h(b,t)&=5+2-6=1,\\
 w_h(c,t)&=2+4-6=0. & & & &
\end{aligned}
\]
Every reduced weight is $\ge 0$, as promised, and the edges on the shortest-path tree
($sb,ba,ac,ct$) have reduced weight $0$ -- they are the ``tight'' edges. (Edge $bt$ has
reduced weight $1$: the route $s\to b\to t$ costs $7$, not the optimal $6$, so $bt$ is
slack, not tight.) Run
Dijkstra on $w_h$: it is now legal, and Lemma~\ref{ch07:reduced-preserves} converts the
answers back via $\dist_w(u,v)=\dist_{w_h}(u,v)-h(u)+h(v)$.
\end{labbox}

The two panels below compare the original negative-weight graph alongside its potential-reweighted counterpart under $h(v)=\dist(s,v)$, where every shortest-path tree edge becomes tight ($w_h=0$) and all remaining edges have positive slack ($w_h>0$).

\begin{center}
\begin{tikzpicture}[
    >=Stealth,
    line width=0.9pt,
    v/.style={circle,draw=nodeblue,fill=blue!6,minimum size=8mm,inner sep=0pt,font=\small},
    vsrc/.style={circle,draw=sourceblue,fill=sourceblue!25,minimum size=8mm,inner sep=0pt,font=\small},
    negedge/.style={draw=trapred!90!black,->,very thick},
    tightedge/.style={draw=targetgreen!85!black,->,very thick},
    slackedge/.style={draw=edgegray,->,dashed},
    normedge/.style={draw=edgegray,->},
    pot/.style={font=\tiny\bfseries,text=structurepurple!90!black},
    paneltitle/.style={font=\footnotesize\scshape,text=gray!85!black}
]
  % Panel (a): Original weights w with potentials h(v)
  \begin{scope}[xshift=0cm]
    \node[paneltitle] at (2.5, 2.35) {(a) Original weights $w$ and potentials $h(v)$};
    \node[vsrc] (s1) at (0,0) {$s$};
    \node[v] (a1) at (2.3,1.25) {$a$};
    \node[v] (b1) at (2.3,-1.25) {$b$};
    \node[v] (c1) at (4.8,1.25) {$c$};
    \node[v] (t1) at (4.8,-1.25) {$t$};

    \node[pot,above=1.5pt of s1] {$h{=}0$};
    \node[pot,above=1.5pt of a1] {$h{=}1$};
    \node[pot,below=1.5pt of b1] {$h{=}2$};
    \node[pot,above=1.5pt of c1] {$h{=}4$};
    \node[pot,below=1.5pt of t1] {$h{=}6$};

    \draw[normedge] (s1) -- node[above left,font=\scriptsize]{$4$} (a1);
    \draw[normedge] (s1) -- node[below left,font=\scriptsize]{$2$} (b1);
    \draw[negedge]  (b1) -- node[left,font=\scriptsize\bfseries,text=trapred!90!black]{$-1$} (a1);
    \draw[normedge] (a1) -- node[above,font=\scriptsize]{$3$} (c1);
    \draw[normedge] (b1) -- node[below,font=\scriptsize]{$5$} (t1);
    \draw[normedge] (c1) -- node[right,font=\scriptsize]{$2$} (t1);
    \draw[normedge] (a1) to[bend left=12] node[above right,font=\scriptsize]{$6$} (t1);
  \end{scope}

  % Panel (b): Reduced weights w_h = w(u,v) + h(u) - h(v) >= 0
  \begin{scope}[xshift=7.2cm]
    \node[paneltitle] at (2.5, 2.35) {(b) Reduced costs $w_h(u,v)=w(u,v)+h(u)-h(v)\ge 0$};
    \node[vsrc] (s2) at (0,0) {$s$};
    \node[v] (a2) at (2.3,1.25) {$a$};
    \node[v] (b2) at (2.3,-1.25) {$b$};
    \node[v] (c2) at (4.8,1.25) {$c$};
    \node[v] (t2) at (4.8,-1.25) {$t$};

    \draw[slackedge] (s2) -- node[above left,font=\scriptsize]{$3$} (a2);
    \draw[tightedge] (s2) -- node[below left,font=\scriptsize\bfseries,text=targetgreen!85!black]{$0$} (b2);
    \draw[tightedge] (b2) -- node[left,font=\scriptsize\bfseries,text=targetgreen!85!black]{$0$} (a2);
    \draw[tightedge] (a2) -- node[above,font=\scriptsize\bfseries,text=targetgreen!85!black]{$0$} (c2);
    \draw[slackedge] (b2) -- node[below,font=\scriptsize]{$1$} (t2);
    \draw[tightedge] (c2) -- node[right,font=\scriptsize\bfseries,text=targetgreen!85!black]{$0$} (t2);
    \draw[slackedge] (a2) to[bend left=12] node[above right,font=\scriptsize]{$1$} (t2);
  \end{scope}
\end{tikzpicture}
\end{center}

\begin{connectionbox}
\textbf{Johnson's algorithm} (all-pairs chapter) applies this reweighting directly: one Bellman--Ford
from an artificial super-source computes $h(v)=\dist(q,v)$ in $\Oh(nm)$, reweights to make
everything nonnegative, then runs Dijkstra from every vertex in $\Oh(n(m+n\log n))$ -- far
better than $n$ Bellman--Fords on sparse graphs. \textbf{Min-cost flow's} successive-shortest-
path algorithm keeps a potential equal to current distances so that each augmenting search is
a nonnegative-weight Dijkstra despite residual edges with negative cost. \textbf{$A^\ast$
search} is Dijkstra with potential $h(v)=$ heuristic estimate of $\dist(v,t)$; an
\emph{admissible, consistent} heuristic is exactly a feasible potential, which is why
consistent $A^\ast$ never needs to re-expand a node -- it is Dijkstra on reduced costs.
\end{connectionbox}

\section{Putting the Pieces Together}
\label{ch07:summary}

\begin{center}
\begin{tabular}{lll}
\toprule
Setting & Algorithm & Time \\
\midrule
Nonnegative weights      & Dijkstra (binary heap)  & $\Oh((n+m)\log n)$ \\
Nonnegative, dense       & Dijkstra (Fibonacci/array) & $\Oh(m+n\log n)$ / $\Oh(n^2)$ \\
Arbitrary weights        & Bellman--Ford           & $\Oh(nm)$ \\
DAG (any weights)        & Topological sweep       & $\Oh(n+m)$ \\
Arbitrary, all-pairs     & Johnson (potentials + Dijkstra) & $\Oh(nm+n^2\log n)$ \\
\bottomrule
\end{tabular}
\end{center}

The chapter rests on three ideas. Every algorithm in it is a schedule of \emph{relaxation}
steps. The \emph{optimality condition} $(\star)$ -- no edge improvable -- is both a
stopping rule and, summed around a cycle, the negative-cycle test. And \emph{potentials}
recast the problem as a shifted potential field, letting fast nonnegative-weight machinery
solve negative-weight problems and tying shortest paths to linear-programming duality and heuristic search.

In Chapter~8, we build on this foundation to explore advanced shortest-path techniques:
goal-directed $A^\ast$ search, bidirectional heuristics, Dial's algorithm for small integer weights,
and scaling methods that push past standard comparison bounds.

\section{Exercises}
\label{ch07:exercises}

\begin{exercise}[Relaxation order]
Give a $4$-vertex weighted DAG and an order of edge relaxations for which a single pass of
Bellman--Ford (one sweep over all edges) does \emph{not} yield correct distances, then give
an order on the same graph for which it does. Explain the difference using
Lemma~\ref{ch07:path-relax}.
\end{exercise}

\begin{exercise}[Dijkstra breaks gracefully?]
Construct a directed graph with a single negative edge on which Dijkstra (settle-once)
returns a \emph{wrong} distance, and a different single-negative-edge graph on which it
happens to return the \emph{correct} distances. What property of the negative edge's
position decides the outcome?
\end{exercise}

\begin{exercise}[Counting shortest paths]
Modify Dijkstra to also compute, for every $v$, the \emph{number} of distinct shortest
paths from $s$ to $v$ (assume nonnegative weights and that this count fits in a machine
word). Prove your update rule correct and argue it adds only $\Oh(m)$ work. Where does the
argument need ``settle once''?
\end{exercise}

\begin{exercise}[$d$-ary tuning]
For a graph with $n=10^4$ vertices and $m=10^6$ edges, compute the heap arity $d$ that
minimizes the asymptotic Dijkstra cost $\Oh\!\left(nd\log_d n + m\log_d n\right)$, and
compare the resulting bound numerically against a binary heap and a Fibonacci heap.
\end{exercise}

\begin{exercise}[Extracting the negative cycle]
Prove that if Bellman--Ford's final sweep finds an improvable edge $(u,v)$, then following
$\pi$ pointers backward from $v$ for $n$ steps lands on a vertex that lies on a negative
cycle. Turn this into an algorithm that prints the cycle in $\Oh(n)$ time.
\end{exercise}

\begin{exercise}[Shortest path with exactly $k$ edges]
Adapt the Bellman--Ford DP to compute, for a fixed $k$, the minimum-weight walk from $s$ to
$t$ using \emph{exactly} $k$ edges (walks may repeat vertices). What is the running time, and
why can you not simply stop relaxing early?
\end{exercise}

\begin{exercise}[Bottleneck shortest paths]
Define the weight of a path as the \emph{maximum} edge weight on it (the bottleneck), and ask
for the path minimizing this maximum. Show that a Dijkstra-like algorithm with the
relaxation $d[v]\gets\min\bigl(d[v],\max(d[u],w(u,v))\bigr)$ solves it, prove correctness,
and state the time bound. How does this connect to minimum spanning trees?
\end{exercise}

\begin{exercise}[Potentials by hand]
For the running graph with $a\to t$ changed to weight $-4$, run one Bellman--Ford from an
added super-source to get a feasible potential $h$, list all reduced weights, verify they are
nonnegative, and confirm that Dijkstra on the reduced graph reproduces the true distances
after the back-conversion of Lemma~\ref{ch07:reduced-preserves}.
\end{exercise}

\begin{exercise}[$A^\ast$ as reduced-cost Dijkstra]
Let $h$ be a \emph{consistent} heuristic estimating distance to the target $t$, meaning
$h(u)\le w(u,v)+h(v)$ for every edge and $h(t)=0$. Show that running Dijkstra with edge
weights $w_h(u,v)=w(u,v)-h(u)+h(v)$ is exactly $A^\ast$ with heuristic $h$, that all reduced
weights are nonnegative, and therefore that consistent $A^\ast$ never re-expands a settled
node.
\end{exercise}

\begin{exercise}[Lazy vs.\ eager heap]
Implement both \textsc{Dijkstra} (with decrease-key) and \textsc{Dijkstra-Lazy} (push-stale)
and construct a family of graphs on which the lazy heap stores $\Th(m)$ entries
simultaneously. Confirm empirically that both run in $\Oh(m\log n)$ but compare their
constant factors and peak memory.
\end{exercise}

\begin{exercise}[Why $n-1$ and not fewer]
Exhibit, for every $n$, a graph on $n$ vertices and an edge-processing order for which
Bellman--Ford genuinely needs all $n-1$ rounds -- the distance of some vertex is still wrong
after $n-2$ rounds. (Hint: a path graph processed back to front.)
\end{exercise}
